\section*{Chapitre \ref{ch:s2:the-central-risk-book} --- Le central risk book}

\begin{solution}{exo:s2:the-central-risk-book:1}
2~millions de dollars croisent~; le net est un achat de 1~million de dollars. Sur les 5~millions de dollars bruts, 80~\% s'annulent par compensation.
\end{solution}

\begin{solution}{exo:s2:the-central-risk-book:2}
$2.326 \times 369\,000 \approx \$858\,000$~; avec l'écart-type non arrondi, 859\,000~\$.
\end{solution}

\begin{solution}{exo:s2:the-central-risk-book:3}
$\$1\text{M} \times (0.0003 + 0.7 \times 0.02 \times \sqrt{1/50}) = \$1\text{M} \times (0.0003 + 0.00198) \approx \$2\,280$.
\end{solution}

\begin{solution}{exo:s2:the-central-risk-book:4}
L'impact par dollar croît avec la taille de la transaction, si bien que les petites transactions coûtent peu par dollar. La compensation supprime les parties opposées des flux, et le net
qui subsiste est négocié en un seul ordre plus grand, avec davantage d'impact par dollar~; l'économie de coût est plus faible que la part du volume croisé (19.5~\%
contre 47.0~\%).
\end{solution}

\begin{solution}{exo:s2:the-central-risk-book:5}
Deux raisons~: les flux clients ultérieurs compensent une partie de ce qui est conservé, si bien qu'il n'a jamais besoin d'être négocié~; ce qui est négocié se fait par tranches quotidiennes plus petites, avec moins d'impact par dollar.
Le prix à payer est le risque résiduel, comme le montre la négociation partielle vers une cible de Garleanu et Pedersen, qui arbitre entre coût et risque.
\end{solution}

\begin{solution}{exo:s2:the-central-risk-book:6}
La composante commune s'additionne sur l'ensemble des desks et ne se compense jamais. Avec davantage de desks, les parties indépendantes se compensent pour un montant égal à $\sqrt D$ fois celui d'un seul desk, tandis que la
partie commune croît comme $D$, si bien qu'elle représente une part plus grande du net, et les transactions nettes sont plus importantes, avec davantage d'impact par dollar.
\end{solution}

\begin{solution}{exo:s2:the-central-risk-book:7}
35~\% par jour~: 89\,300~\$ par jour, avec une valeur en risque de 449\,000~\$. Une limite plus stricte achète moins de risque, à un coût deux fois supérieur à celui de la limite de 1~million de dollars.
\end{solution}

\begin{solution}{exo:s2:the-central-risk-book:8}
À 5~\% par jour, le P\&L porté a un écart-type de 684\,000~\$ et une valeur en risque à un jour, à 99~\%, d'environ 1.59~million de dollars, bien au-dessus de la limite de 1
million de dollars. Le coût tombe à 20\,300~\$ par jour, soit 93.9~\% de moins que desk par desk, parce que le portefeuille porte bien davantage de risque~; l'économie est le prix d'un risque,
non une amélioration gratuite.
\end{solution}

\begin{solution}{pb:s2:the-central-risk-book:1}
\begin{enumerate}
\item Un portefeuille qui reprend les risques que les desks acquièrent auprès des clients, les compense et gère le net~; la \omterm{def:s2:the-central-risk-book:pooling}{mutualisation des risques} combine les positions avant la couverture, de sorte que les positions opposées s'annulent.
\item Cinq desks, 100 actions, 500 jours, écart-type des flux de 0.5~million de dollars par nom et par jour, avec 10~\% de variance commune~; un modèle factoriel à trois facteurs avec une
volatilité idiosyncratique de 1.5~\%~; une demi-fourchette de 3~points de base et un impact en racine carrée de coefficient 0.7.
\item 335\,000~\$ par jour, desk par desk, 270\,000~\$ mutualisés (19.5~\% de moins)~; 47.0~\% des flux bruts sont croisés.
\item Internaliser le risque d'un client en anticipation d'un flux de compensation, ou l'externaliser en le couvrant immédiatement.
\item Avec des flux indépendants, l'économie monte à 48.0~\% avec dix desks~; avec une partie commune, elle culmine à 19.7~\% avec six desks et tombe à 17.1~\% avec dix.
\item Elle ne se compense pas, et croît plus vite que la part compensée.
\item L'impact temporaire suit une puissance 3/5 du taux de transaction plutôt qu'une racine carrée~; dans les deux cas, l'impact croît plus vite que la taille.
\item L'impact par dollar est plus élevé sur la transaction nette, plus grande, que sur les petites fractions opposées.
\item Couvrir immédiatement les facteurs communs avec des instruments liquides et laisser le reste à négocier ou à compenser.
\item Négocier chaque jour une part fixe de l'inventaire avec des contrats à terme sur le reste pondéré par les bêtas~: de 269\,700~\$ sans aucun risque à 100~\%, à 20\,300~\$ par jour et
un écart-type du P\&L quotidien de 684\,000~\$ à 5~\%.
\item Elle supprime environ un cinquième du risque porté pour environ 500~\$ par jour.
\item 15~\%~: 45\,300~\$ par jour et une valeur en risque de 859\,000~\$.
\item La mutualisation économise 19.5~\% du coût de couverture~; une couverture patiente dans une limite de valeur en risque de 1~million de dollars économise 86.5~\%, en portant un
écart-type quotidien du P\&L de 369\,000~\$.
\item Que la couverture de chaque desk soit exécutée comme si ce desk était seul sur le marché~; en réalité, les ordres des desks se rencontrent sur le même marché.
\item Facturer aux desks le risque qu'ils transmettent et les créditer de la compensation qu'ils apportent, de sorte qu'ils ne déversent pas leur risque et ne contournent pas le portefeuille.
\item Les flux de compensation cessent d'arriver, l'inventaire croît, et le taux doit monter quand la négociation est la plus chère.
\item L'information des clients d'un desk et la négociation d'un autre~; les limites du portefeuille et les incitations des desks.
\item \omterm{def:s2:the-central-risk-book:residual}{Couverture du risque résiduel} par facteurs.
\item Le desk du chapitre~24 internalise un seul flux client~; le \omterm{def:s2:the-central-risk-book:crb}{central risk book} internalise à travers les desks.
\item Compensation et patience~: le coût que la banque ne paie pas pour couvrir des risques qui se compensent ou peuvent attendre.
\end{enumerate}
\end{solution}

\begin{solution}{iq:s2:the-central-risk-book:1}
Un portefeuille qui recueille le risque des clients de plusieurs desks, le compense et ne couvre que le reliquat, d'abord les risques de facteurs, puis le reste patiemment. Il économise les
fourchettes et l'impact de la couverture de risques qui se compensent, et donne à la banque une vue unique de ses expositions.
\end{solution}

\begin{solution}{iq:s2:the-central-risk-book:2}
Tracer la frontière coût-risque sur les flux, coûts et volatilités effectifs du portefeuille, puis choisir le point qui minimise le coût plus un prix du risque, ou le
point le moins cher dans la limite de risque~; reprendre le calcul lorsque les flux ou la volatilité changent.
\end{solution}

\begin{solution}{iq:s2:the-central-risk-book:3}
Au prix milieu au moment du croisement, chaque desk étant facturé de la même manière, et une charge de transfert reflétant le coût qu'il aurait payé à l'extérieur
diminué d'une part de l'économie~; jamais à un prix qui favorise les clients d'un desk au détriment de ceux de l'autre.
\end{solution}

\begin{solution}{iq:s2:the-central-risk-book:4}
Valeur en risque (VaR) et perte de stress~; expositions aux facteurs (marché, secteur, style)~; concentration par nom au regard de sa liquidité~; délai de liquidation~; limites sur
l'ancienneté de l'inventaire~; et escalade lorsque le flux de compensation se tarit.
\end{solution}

\begin{solution}{iq:s2:the-central-risk-book:5}
Les positions et transactions de chaque desk, par instrument et dans le temps, diffusées avec des données de référence cohérentes et des horodatages~; compensation par sous-jacent~; expositions aux
facteurs calculées centralement~; rapprochement avec les portefeuilles des desks~; et une piste d'audit des croisements.
\end{solution}

\begin{solution}{iq:s2:the-central-risk-book:6}
Desk par desk, le coût linéaire espéré vaut $D c\,E|q| = D c s\sqrt{2/\pi}$~; mutualisé, l'écart-type du net vaut $\sqrt D s$, donc $c\sqrt D s\sqrt{2/\pi}$~:
soit un facteur $\sqrt D$ de moins. Le coût d'impact est $k E|q|^{3/2} \propto s^{3/2}$ par desk, soit $D$ au total~; mutualisé $\propto (\sqrt D s)^{3/2} = D^{3/4}s^{3/2}$~:
soit un facteur $D/D^{3/4} = D^{1/4}$ de moins.
\end{solution}
